\documentclass[10pt,a4paper]{article}

\usepackage[margin=1.05in]{geometry}
\usepackage{setspace}
\usepackage{amsmath,amssymb,amsthm}
\usepackage{booktabs}
\usepackage{array}
\usepackage{caption}
\usepackage{enumitem}
\usepackage{titlesec}
\usepackage{microtype}
\usepackage{hyperref}
\usepackage{fancyhdr}

\hypersetup{colorlinks=true,linkcolor=black,urlcolor=black,citecolor=black}
\captionsetup{font=small,labelfont=bf,skip=4pt,justification=raggedright}

\titleformat{\section}{\normalsize\bfseries}{}{0em}{}
\titleformat{\subsection}{\small\bfseries}{}{0em}{}
\titlespacing*{\section}{0pt}{1.35em}{0.5em}
\titlespacing*{\subsection}{0pt}{0.95em}{0.35em}

\theoremstyle{plain}
\newtheorem{proposition}{Proposition}
\theoremstyle{definition}
\newtheorem{definition}{Definition}
\theoremstyle{remark}
\newtheorem{observation}{Observation}

\setstretch{1.14}
\setlist[itemize]{leftmargin=1.15em,itemsep=0.1em,topsep=0.2em}

\pagestyle{fancy}
\fancyhf{}
\fancyhead[L]{\small\itshape Elements of Position}
\fancyhead[R]{\small\itshape The Geometric Kernel}
\fancyfoot[C]{\small\thepage}
\renewcommand{\headrulewidth}{0pt}

\begin{document}

\begin{center}
\vspace*{2.2em}
{\large Elements of Position}\\[1.6em]
{\LARGE\bfseries The Geometric Kernel}\\[0.6em]
{\small minimal run --- 18 September 2026}\\[1.4em]
{\normalsize Justin Erholtz}\\[0.35em]
{\small 18 September 2026}
\end{center}

\vspace{1.6em}

\begin{center}
{\large What the walk measures}
\end{center}

\vspace{1.4em}

\noindent
Paper~I named the closed unit and the pair.
This paper is the walk that produces those measurements from two operators, a station, and a count.
Nothing else is in the operators.
What is named after the fact is read off the tables the walk writes.

The source is \texttt{EOP\_Geometric\_Kernel\_minimal.py}.
The tables are \texttt{EOP\_crossings\_minimal.csv}, \texttt{EOP\_pairs\_minimal.csv}, \texttt{EOP\_summary\_minimal.csv}, from a run with shear \(h=0.002\) and \(200\,000\) steps (\(127\) horizon hits).
Those numbers are a sheet size and a depth of count.
The identities do not belong to that sheet.

\vspace{1.2em}

\noindent
\textit{A note.}
Claims that can be proved from the operators are proved as identities.
Claims that are counts are reported from the tables.
The two are kept apart.
Trigonometry does not enter the step.
It may appear only as a named comparison after the tables exist.
Named comparisons sit \emph{between} the two integers and the two bands the walk actually writes.
They are not values the last row equals.

\newpage
\tableofcontents
\newpage

\section{What is in the operators}

Start on the diameter: \((1,0)\).
A heading is a pair of coordinates whose length after the bind is \(1\).

\begin{definition}[Shear]
For a kick \(h>0\),
\[
T_h(x,y)=(x-hy,\,y+hx).
\]
\end{definition}

\begin{definition}[Bind]
\[
N(x,y)=\frac{(x,y)}{\sqrt{x^2+y^2}}.
\]
\end{definition}

\begin{definition}[Step]
\[
\mathrm{step}=N\circ T_h.
\]
\end{definition}

No angle is an argument of \(T_h\) or of \(N\).
The shear is an equal opposite mix of the two coordinates.
The bind restores the quadratic \(x^2+y^2=1\).

If the last heading already had length \(1\), the unbound length after the shear is an identity of the sheet:
\[
\bigl|T_h(u)\bigr|=\sqrt{1+h^2}.
\]
The bind leftover of one kick is therefore \(\sqrt{1+h^2}-1\), constant in the step index.
On the recorded run this is \(1.999998\times 10^{-6}\).
The Euclidean distance between two successive bound headings is a function of \(h\) only.
On the recorded run it is \(0.001999997\).

Those two numbers describe the grain of the sheet.
They are not a position on the pair.

The shape of \(T_h\) is not empty.
Equal opposite off-diagonals are Euclid's quarter-turn generator.
The bind is Euclid's length.
Drop either and the walk below does not measure what it measures.
Neither operator contains a finished turn, a named half-length, or the inverse pair.

\section{The station and the horizon}

The heading begins on \(y=0\).
A hit is that heading crossing \(y=0\) again.

\begin{definition}[First gap]
\(H_0\) is the step index of the first hit.
\end{definition}

On the recorded run \(H_0=1571\), side down, heading
\[
(-0.99999992,\,-0.00040316).
\]
The heading is on the far wall, one kick of vertical.

Stand on the line you count and the first event is the far side of that line.
Stand off the line and the first gap is only how far the station was from the wall.
This run stands on the diameter.

\begin{observation}
The unit heading from the start to the first far wall has length \(2\).
The unit heading from the first wall to the second has length \(2\).
Hits to restore the starting side: \(2\).
That \(2\) is not a coefficient in the pair formula.
\end{observation}

On the recorded run both wall chords are \(1.99999996\).

\section{The count is two integers}

After the first hit the tables show only two gap values, \(1570\) and \(1571\).
On \(127\) hits:
\[
1571\text{ occurs }102\text{ times},\qquad
1570\text{ occurs }25\text{ times}.
\]
The first hit is the first gap \(H_0=1571\).
The remaining word of long and short gaps is
\[
\texttt{LLLSLLLLS}\ldots\texttt{LLLLSLL},
\]
twenty-five copies of \(\texttt{LLLLS}\) and then \(\texttt{LL}\) because the step budget ended mid-word.
That \(5\) is the period of this sheet's two-integer payment.
It is not the pentagon of Paper~III.

Mean gap \(1570.80315\).
Mean gap times the kick: \(3.14160630\).

\begin{observation}
The product of the mean gap and the kick is a reading of the sheet.
It is not inserted into \(T_h\) or \(N\).
The discrete walk cannot land on one length, so it writes both neighbors.
\end{observation}

A named comparison after the tables exist: \(\pi=3.14159265\), and \(\pi/\operatorname{atan}(h)=1570.798421\).
Those names sit between \(1570\) and \(1571\).
\(H_0\cdot h=3.142\) is the first gap only --- one long step, coarser than the mean.

Frozen scale \(j/H_0\) uses the first gap as a ruler.
Snap scale is the hit index \(m\).
They agree at \(m=1\).
They part by one first-gap unit per short gap.
At \(m=127\),
\[
m-\frac{j}{H_0}=0.01591343=\frac{25}{1571}.
\]
The frozen drift \emph{is} the short count over \(H_0\).

\section{The heading ratchets and resets}

On a long hit the heading is off the axis by about a kick, then two, then three, then four.
On the next short hit it is put back.

\begin{observation}
On every short hit of the recorded run,
\[
\frac{|h_y|}{m}=\alpha,\qquad
\alpha=3.15762023\times 10^{-6}.
\]
The same \(\alpha\), to the digits of the mean, from \(m=5\) to \(m=125\).
\end{observation}

Short hit: leftover proportional to how far the count has run.
Long hits: ratchet off the wall.
That is snap versus frozen in the heading, not as a speech.
Snap to the event and the heading stays a horizon hit.
Freeze the first gap and the heading walks off the axis until a short gap pays the accumulated leftover.

\section{The pair}

After the walk, bind a part and its inverse to the heading at hit \(m\):
\[
r=m,\qquad \frac1r=\frac1m.
\]

\begin{proposition}
For every \(m\ge 1\),
\[
r\cdot\frac1r=1,\qquad
\mathrm{GM}\Bigl(r,\tfrac1r\Bigr)=1.
\]
\end{proposition}

\begin{proof}
Immediate.
\end{proof}

The crossings table records \(\mathrm{gm}=1\) on every row, to machine precision.

\begin{definition}[Leftover of the pair]
\[
\tau(m)=\mathrm{AM}-\mathrm{GM}
=\frac{m+1/m}{2}-1.
\]
\end{definition}

\begin{proposition}
\(\tau(m)=0\) if and only if \(m=1\) (for \(m>0\)).
\end{proposition}

\begin{proof}
\(m+1/m=2\) multiplies to \((m-1)^2=0\).
\end{proof}

\begin{proposition}
\(\tau(m)=\tau(1/m)\) for every \(m>0\).
\end{proposition}

\begin{proof}
Substitute.
\end{proof}

Hit \(1\): pair \((1,1)\), \(\tau=0\).
The lock.
The first horizon.

Hit \(2\): pair \((2,1/2)\), \(\mathrm{AM}=5/4\), \(\tau=1/4\).
Both horizons.
The dated midpoint of Paper~I, now that both ends exist.

The involution is \(m\leftrightarrow 1/m\).
\texttt{EOP\_pairs\_minimal.csv} is that formula on each row.
It is not a table of integer partners.
The only integer partner in the count is \(m=1\) with itself.
There is no \(4/k\).
That \(4\) belonged to an earlier units choice \(r=m/2\), which moved the kiss to hit \(2\) and wrote \(\{2,1/2\}\) on hit \(1\).
Same pairs, wrong clock.

\section{Chord and path: two bands}

Place the outward wall on the heading:
\[
P_m=m\,(h_x,h_y).
\]
The straight line from \(P_{m-1}\) to \(P_m\) is the leg chord.
The summed displacement of \(P\) along the bound headings of that leg, radius linear from \(m-1\) to \(m\) in the leg's own gap, is the leg path.
No kick is treated as an angle in that sum.

\begin{proposition}
If the heading at hit \(m-1\) is the opposite of the heading at hit \(m\), the chord is
\[
(m-1)+m=2m-1.
\]
\end{proposition}

\begin{proof}
Opposite unit headings \(u\) and \(-u\):
\[
\bigl|m(-u)-(m-1)u\bigr|=(m+m-1)|u|=2m-1.
\]
\end{proof}

The table: \(m=2\) chord \(3.000000\); \(m=3\) chord \(5.000000\); \(m=4\) chord \(7.000000\).
The short-gap rows miss \(2m-1\) by a few parts in \(10^{6}\), the same leftover that shows in \(h_y\).

Chord over path does not climb to a last-row value of \(2/\pi\).
It writes two bands.

\begin{center}
\begin{tabular}{@{}lccc@{}}
\toprule
 & last short & last long & \(2/\pi\) \\
\midrule
chord/path & \(0.63694134\) & \(0.63653616\) & \(0.63661977\) \\
\bottomrule
\end{tabular}
\end{center}

Short legs sit above the named comparison.
Long legs sit below it.
The last row of this run is a long hit, so the last chord/path \(0.636536\) is the lower band, not a leftover angle treated as a second constant.
Mean of all legs after \(m=1\): \(0.636381\).

The named comparison sits between the bands, as \(\pi/\operatorname{atan}(h)\) sits between the two gap integers.

\section{The leftover integrated}

From the kiss at \(m=1\) outward,
\[
\int_1^N\tau(k)\,dk
=\frac{N^2-1}{4}+\frac12\ln N-(N-1).
\]
On the recorded run \(N=127\):
\[
\text{closed}=3908.422094.
\]
The sum of \(\tau(m)\) on the integer hits is \(3939.712667\).
The trapezoid correction is
\[
\int\tau+\frac12\bigl(\tau(1)+\tau(N)\bigr)=3939.674062.
\]
Sum minus trapezoid: \(0.038605\).
That difference is the next Euler--Maclaurin term on a function already named.
It is not a third leftover.

\section{What the three files hold}

\subsection{\texttt{EOP\_summary\_minimal.csv}}

Sheet: \(h\), leftover of one kick, heading chord, \(H_0\).
Count: two gap integers, the word of longs and shorts, mean gap times kick.
Heading: \(\alpha=|h_y|/m\) on short hits, frozen drift \(=n_{\mathrm{short}}/H_0\).
Pair: kiss \((1,1)\) at \(m=1\), pair \((2,1/2)\) at \(m=2\).
Bands: last short chord/path, last long chord/path.
Wedge: sum, integral, trapezoid, their difference.
Comparisons after the tables: \(\pi\), \(\pi/\operatorname{atan}(h)\), \(2/\pi\).

\subsection{\texttt{EOP\_crossings\_minimal.csv}}

One row per hit.
Forced on every row: \(\mathrm{gm}=1\), \(r\cdot(1/r)=1\), \(\tau=(m+1/m)/2-1\), heading on a horizon.
\texttt{gap\_kind} is \texttt{first}, \texttt{long}, or \texttt{short}.
\texttt{hy\_over\_m} is constant on the short rows.

\subsection{\texttt{EOP\_pairs\_minimal.csv}}

The formula \(m\), \(m\), \(1/m\), AM, GM, \(\tau\).
No empty partner columns.

\section{What is not a reading of this walk}

\(\Phi\) is not among the readings.
Process~B remains a second street, walked after arrival, as in Paper~II.

The period-\(5\) gap word on this sheet is the two-integer payment for \(\pi/\operatorname{atan}(h)\approx 1570.798\).
It is not a docking with Paper~III.

An unequal mix \((x-ay,\,y+bx)\) still closes, but mean times kick tracks \(\pi h/\sqrt{ab}\), not the square-mix reading.
A one-sided mix from this station does not close.
The equal opposite shear is a choice of kind.
It is Euclid's \(90^\circ\) in the operator.
It is not \(\pi\) in the operator.

\section{What stands}

The operators produce a heading that counts crossings of the diameter it stood on.

The count is two neighboring integers.
Mean gap times kick sits between them.

The heading ratchets off the wall on long gaps and is reset on short gaps to \(|h_y|=m\alpha\).

Scale is the hit index.
The pair at that scale has product \(1\) and leftover \((m+1/m)/2-1\), zero only at \(1\).
At two hits the pair is \(\{2,1/2\}\).
The unit heading already had wall chords of length \(2\).

Chord over the actual path of the outward wall writes two bands.
The named comparison \(2/\pi\) sits between them.

Named after the measure, not before it, and not as a last-row equality:
mean gap times kick; wall chord \(2\); two bands around \(2/\pi\).
None of those names enters \(T_h\) or \(N\).

What is not granted: a last decimal as closure of the heading; a completed run of the count; \(\Phi\) as an output of this walk; a leftover of the pair treated as a capacity or a field; the extra \(2\) in an older radius \(m/2\); a partner table built by rounding \(4/k\).

The first sentence of Paper~I is still the sentence.
The kernel is how the walk answers it.

\end{document}
